\section{给定预算，参数与数据怎么分}
\label{sec:14-Deep-Learning/09-Scaling-and-Adaptation/02}

预算批下来了：一个固定的计算量 \(C\)，用完为止。现在要填两个数字——模型做多大，训多少 token。这两个数字不独立，因为算力被它们的乘积吃掉：参数翻倍，同样的算力就只够跑一半的数据。

一个流行过很多年的答案是``参数优先''：先把模型堆到显存的上限，剩下多少算力就训多少数据。这个答案不是从任何计算里得出的，它来自一种朴素印象——模型越大越好。上一节那条幂律恰好给了检验它的工具，因为把两条幂律和一个预算约束写在一起，就是一个标准的等式约束优化问题，而它有解析解。

\subsection{预算约束长什么样}

训练一次的浮点运算量近似是

\[
  C \approx 6 N D,
\]

\(N\) 是参数量，\(D\) 是训练 token 数。常数 \(6\) 值得拆开看，它不是拟合来的。前向传播中每个参数对每个 token 大致参与一次乘法和一次加法，是 \(2\) 次浮点运算；反向传播要算两组导数——对激活的和对权重的——各花掉与前向相当的一份，合计约 \(4\) 次。加起来每个参数每个 token 约 \(6\) 次。这个计数忽略了注意力矩阵那部分与序列长度平方相关的开销（见\hyperref[sec:14-Deep-Learning/05-Attention-and-Sequence-Models/01]{``Self-Attention 是内容决定的加权读取''}），在上下文不太长时它占比不大，上下文拉长后这条近似会明显偏低。

损失那边，把上一节两条幂律并起来：

\[
  L(N, D) = L_\infty + \frac{A}{N^{\alpha}} + \frac{B}{D^{\beta}},
  \qquad A, B > 0,\quad \alpha, \beta > 0 .
\]

这个可加形式本身也是经验假设——它断言参数不足和数据不足这两种缺陷互不纠缠，各自贡献一项。真实情况下两者当然会相互影响，但在拟合区间内这个形式够用，而且它让接下来的计算变得干净。

\subsection{一次 Lagrange 乘子}

在 \(6ND = C\) 的约束下最小化 \(L\)。常数项 \(L_\infty\) 不影响最优点，去掉。写出 Lagrangian

\[
  \mathcal{L}(N, D, \lambda) = \frac{A}{N^{\alpha}} + \frac{B}{D^{\beta}} + \lambda\,(6ND - C).
\]

对两个变量求偏导并令其为零：

\[
  \begin{aligned}
    \frac{\partial \mathcal{L}}{\partial N} &= -\alpha A N^{-\alpha-1} + 6\lambda D = 0, \\[2pt]
    \frac{\partial \mathcal{L}}{\partial D} &= -\beta B D^{-\beta-1} + 6\lambda N = 0 .
  \end{aligned}
\]

第一式两边乘 \(N\)，第二式两边乘 \(D\)，右侧都出现 \(6\lambda ND = \lambda C\)：

\[
  \alpha A N^{-\alpha} = \lambda C = \beta B D^{-\beta}.
\]

这一步是整个推导的支点。它说的是：最优处，参数不足带来的损失与数据不足带来的损失，各自按指数加权之后必须相等。哪一边小了，把资源从另一边挪过来就能降损失，所以不可能是最优。

由此 \(N^{\alpha} \propto D^{\beta}\)，即 \(D \propto N^{\alpha/\beta}\)。代回约束 \(ND = C/6\)：

\[
  N \cdot N^{\alpha/\beta} \propto C
  \quad\Longrightarrow\quad
  N^{(\alpha+\beta)/\beta} \propto C
  \quad\Longrightarrow\quad
  N^\star \propto C^{\beta/(\alpha+\beta)},
  \qquad
  D^\star \propto C^{\alpha/(\alpha+\beta)} .
\]

两个指数加起来正好是 \(1\)，与 \(C \approx 6ND\) 相容，这是一处可以自查的地方。

\begin{remark}
条件对账。上面用的是等式约束下的一阶必要条件，它本身只保证驻点。要它同时是全局最小，需要三件事。第一，\(A, B > 0\) 且 \(\alpha, \beta > 0\)，否则损失关于规模不是递减的，问题本身就变了。第二，最优点落在可行域内部而非边界——现实中总有 \(N \ge N_{\min}\) 这类工程下界，若最优解撞上它，KKT 条件里的互补松弛项就不能忽略，这一层的处理见\hyperref[sec:09-Convex-Optimization/05-Duality/03]{``KKT 把最优性写成一组方程''}。第三，唯一性：换成对数变量 \(n = \log N\)、\(d = \log D\)，目标变成 \(A e^{-\alpha n} + B e^{-\beta d}\)，两项都是严格凸的指数函数，约束变成线性的 \(n + d = \log(C/6)\)，于是限制在约束上的目标严格凸，驻点唯一且是全局最小。这三条都成立时，上面的解就是答案。

不构成条件的是幂律本身。\(\alpha\)、\(\beta\)、\(A\)、\(B\) 全部来自拟合，换一个拟合区间、换一种学习率调度，估计值会漂移，而结论里的指数 \(\beta/(\alpha+\beta)\) 对这种漂移是敏感的。数学部分是严格的，输入进去的数不是。
\end{remark}

\subsection{于是应该同比例增长}

结论的形状取决于 \(\alpha\) 和 \(\beta\) 的相对大小，而经验上这两个指数相当接近。取 \(\alpha \approx \beta\)，上式退化成

\[
  N^\star \propto C^{1/2}, \qquad D^\star \propto C^{1/2}.
\]

算力翻十倍，参数量和训练 token 数各自乘以 \(\sqrt{10} \approx 3.2\)——同比例放大，谁也不该被优先。

对照一下``参数优先''的做法：那相当于让 \(N \propto C\)、\(D\) 基本不动。放进上面的最优性条件看，此时 \(\alpha A N^{-\alpha}\) 被压得很小，而 \(\beta B D^{-\beta}\) 原地不动，两边差得越来越远。等式不成立的方向明确地告诉你该怎么改：把算力从参数挪到数据。这正是那批早期大模型的处境——它们不是训坏了，是配比错了；同样的电费本可以换到明显更低的损失。

值得停一下的是这个结论的来源。它不来自新的实验，也不来自架构上的洞见，就是把两条已有的经验曲线和一个预算约束摆在一起求了个极值。工具是第 09 部分的 Lagrange 乘子，翻译的过程是\hyperref[sec:09-Convex-Optimization/01-Optimization-Foundations/01]{``把现实选择翻译成优化问题''}讲的那件事。数学在这里的作用不是提供新信息，而是把已经握在手里的信息里的矛盾逼出来。

\subsection{切点与它移动的方向}

\begin{center}
\begin{tikzpicture}[scale=1.0]

\draw[->,>=stealth] (0.3,0.3) -- (8.5,0.3);
\draw[->,>=stealth] (0.3,0.3) -- (0.3,5.6);
\node[anchor=north] at (4.4,0.12) {\footnotesize 参数量 \(N\)};
\node[anchor=south,rotate=90] at (-0.05,2.9) {\footnotesize 数据量 \(D\)};

% 等算力曲线 ND = C（三个预算）
\draw[dashed] plot[domain=0.335:4.20,samples=90] (\x,{1.69/\x});
\draw[dashed] plot[domain=0.870:6.00,samples=90] (\x,{4.41/\x});
\draw[dashed] plot[domain=1.780:7.80,samples=90] (\x,{9.00/\x});

% 等损失曲线：(N-n_0)(D-n_0)=k，在对角线上与等算力曲线相切
\draw[very thick] plot[domain=0.700:4.20,samples=90] (\x,{0.585+0.5112/(\x-0.585)});
\draw[very thick] plot[domain=1.270:6.00,samples=90] (\x,{0.945+1.3340/(\x-0.945)});
\draw[very thick] plot[domain=2.080:7.80,samples=90] (\x,{1.350+2.7225/(\x-1.350)});

% 切点与最优射线
\draw[dotted,thick] (0.55,0.55) -- (3.45,3.45);
\fill (1.30,1.30) circle (2.2pt);
\fill (2.10,2.10) circle (2.2pt);
\fill (3.00,3.00) circle (2.2pt);
\node[anchor=west] at (3.55,3.50) {\scriptsize 切点连线：\(N\) 与 \(D\) 同比例};

% 计入推理成本后切点沿等算力曲线向左上移动
\draw[<-,>=stealth,very thick] plot[domain=1.45:2.05,samples=40] (\x,{4.41/\x});

% 图例（放在右上角的空白区）
\draw[dashed] (3.70,5.15) -- (4.35,5.15);
\node[anchor=west] at (4.45,5.15) {\scriptsize 等算力曲线 \(ND = C\)};
\draw[very thick] (3.70,4.65) -- (4.35,4.65);
\node[anchor=west] at (4.45,4.65) {\scriptsize 等损失曲线（越往右上损失越低）};
\draw[->,>=stealth,very thick] (3.70,4.15) -- (4.35,4.15);
\node[anchor=west] at (4.45,4.15) {\scriptsize 计入推理成本后切点的移动方向};

\end{tikzpicture}
\end{center}

\emph{虚线是等算力曲线 \(ND = C\)（双曲线，越往右上算力越大），粗实线是等损失曲线（越往右上损失越低）。给定预算求最优，就是找那条虚线上损失最低的点，也就是两族曲线相切的地方；三个切点连成的点线接近对角方向，这就是 \(N^\star\) 与 \(D^\star\) 同比例增长的几何含义。要记住的是右上那个箭头：目标函数里一旦加进推理成本，切点就沿它移动——更小的模型、更多的数据。最优点是目标定义出来的，不是模型本身的属性。}

\subsection{这个最优是对训练成本而言的}

上面整个推导里，``成本''只有一项：训练那一次的 \(6ND\)。模型训完之后要被使用，而使用也要花钱，这笔钱的结构完全不同。

推理一个 token 的开销大致是 \(2N\) 次浮点运算——只有前向，没有反向，所以是 \(2\) 而不是 \(6\)。关键在于它乘的不是训练数据量，而是这个模型一生要服务的 token 总数 \(T\)。总账变成

\[
  C_{\text{total}} \approx 6ND + 2NT .
\]

\(T\) 由部署规模和服务时长决定，与训练无关，而且对一个用得广、活得久的模型，\(2NT\) 可以轻易超过 \(6ND\)。此时该问的问题也变了：不再是``给定 \(C\) 最小化 \(L\)''，而是``给定目标损失 \(L_0\) 最小化 \(C_{\text{total}}\)''。

新问题的解会明显偏向小模型。原因直接写在式子里：\(N\) 在两项里都出现，\(D\) 只在一项里出现，所以增大 \(N\) 的边际代价被 \(T\) 放大，而增大 \(D\) 的开销是一次性的。把一个较小的模型训得远超训练最优的 token 数，从训练账本上看是浪费，从总账上看可能便宜得多。这正是许多面向部署的模型选择``小模型、超长训练''的原因，它不是妥协，是换了目标函数之后的最优解。

这件事在本书里出现过很多次，形式各不相同，判断是同一条：最优点是目标函数的性质，不是被优化对象的性质。改一个字的目标，最优解可以整体挪位。看到``最优配比''这类说法时，第一个要问的不是数值对不对，而是它在对哪个目标最优。

\medskip

无论最优配比怎么算，结论都是同一个方向：想要更低的损失，就得付更多的钱，而且是按 \(2^{1/\alpha}\) 那个倍数付。这条路走到某处必然停下来。于是问题从``怎样把模型做得更好''转向另一半——已经有了一个训好的大模型，怎样在不重新训练它的前提下改变它的行为。下一节从最省的那个办法开始：只学一个低秩的增量。
